\chapter{Natural Transformations}
\label{ch:nattrans}

\section{Definition}

Natural transformations are morphisms between functors. They provide the precise meaning of ``a natural map'' between constructions—a notion that Eilenberg and Mac Lane introduced category theory specifically to capture.

\begin{definition}[Natural transformation]
Let $F, G: \cat{C} \to \cat{D}$ be functors. A \textbf{natural transformation} $\alpha: F \Rightarrow G$ consists of:
\begin{itemize}
  \item For each object $a \in \cat{C}$, a morphism $\alpha_a: Fa \to Ga$ in $\cat{D}$, called the \textbf{component} of $\alpha$ at $a$.
\end{itemize}
These components must satisfy the \textbf{naturality condition}: for every morphism $f: a \to b$ in $\cat{C}$, the following square commutes in $\cat{D}$:
\[
  \begin{tikzcd}
    Fa \arrow[r, "\alpha_a"] \arrow[d, "Ff"'] & Ga \arrow[d, "Gf"] \\
    Fb \arrow[r, "\alpha_b"'] & Gb
  \end{tikzcd}
\]
That is, $Gf \circ \alpha_a = \alpha_b \circ Ff$.
\end{definition}

\begin{remark}
The naturality condition says that the components of $\alpha$ ``commute with the action of morphisms.'' A natural transformation does not just require a map $Fa \to Ga$ for each $a$; it requires this collection of maps to be coherent with the functorial structure.
\end{remark}

\begin{example}[Double dual]
Let $F = \Id_{\Vect_k}$ and $G = (-)^{**}$ (double dual). For each vector space $V$, define $\alpha_V: V \to V^{**}$ by $\alpha_V(v)(\phi) = \phi(v)$ (evaluation). This is a natural transformation $\Id \Rightarrow (-)^{**}$: for any linear map $T: V \to W$, the diagram
\[
  \begin{tikzcd}
    V \arrow[r, "\alpha_V"] \arrow[d, "T"'] & V^{**} \arrow[d, "T^{**}"] \\
    W \arrow[r, "\alpha_W"'] & W^{**}
  \end{tikzcd}
\]
commutes. (When $\dim V < \infty$, $\alpha_V$ is an isomorphism, but the \emph{single} dual $V \to V^*$ cannot be made natural with respect to all linear maps—it requires a choice of basis.)
\end{example}

\begin{example}[Determinant]
The determinant $\det: \mathbf{GL}_n \Rightarrow (-)^\times$ is a natural transformation from the functor sending a commutative ring $R$ to $\mathbf{GL}_n(R)$ (invertible $n \times n$ matrices) to the functor sending $R$ to its unit group $R^\times$.
\end{example}

\begin{example}[Abelianization]
Let $\mathrm{ab}: \Grp \to \Ab$ be the abelianization functor $G \mapsto G/[G,G]$. There is a natural transformation $\eta: \Id_{\Grp} \Rightarrow i \circ \mathrm{ab}$ (where $i: \Ab \hookrightarrow \Grp$ is inclusion), given by the quotient maps $\eta_G: G \to G/[G,G]$.
\end{example}

\begin{definition}[Natural isomorphism]
A natural transformation $\alpha: F \Rightarrow G$ is a \textbf{natural isomorphism} (written $\alpha: F \xrightarrow{\sim} G$ or $F \cong G$) if every component $\alpha_a: Fa \to Ga$ is an isomorphism in $\cat{D}$.
\end{definition}

\begin{lemma}
$\alpha: F \Rightarrow G$ is a natural isomorphism if and only if $\alpha$ is an isomorphism in the functor category $[\cat{C}, \cat{D}]$ (defined below).
\end{lemma}

\section{The Functor Category}

\begin{definition}[Functor category]
Let $\cat{C}$ and $\cat{D}$ be categories. The \textbf{functor category} $[\cat{C}, \cat{D}]$ (also written $\cat{D}^{\cat{C}}$ or $\Fun(\cat{C},\cat{D})$) has:
\begin{itemize}
  \item Objects: functors $F: \cat{C} \to \cat{D}$.
  \item Morphisms: natural transformations $\alpha: F \Rightarrow G$.
  \item Composition: vertical composition (defined below).
  \item Identities: the identity natural transformation $\id_F$, with components $(\id_F)_a = \id_{Fa}$.
\end{itemize}
\end{definition}

\begin{remark}
When $\cat{C}$ is small and $\cat{D}$ is locally small, the functor category $[\cat{C},\cat{D}]$ is locally small. If $\cat{C}$ is not small, there may be set-theoretic issues.
\end{remark}

\begin{example}[Presheaf category]
For a small category $\cat{C}$, the \textbf{presheaf category} $\widehat{\cat{C}} = [\cat{C}\op, \Set]$ is a functor category of fundamental importance. Its objects are contravariant functors $\cat{C}\op \to \Set$ (called \emph{presheaves} on $\cat{C}$), and its morphisms are natural transformations between them. The Yoneda embedding (Corollary~\ref{cor:yoneda-embedding}) embeds $\cat{C}$ fully faithfully into $\widehat{\cat{C}}$.
\end{example}

\section{Vertical Composition}

\begin{definition}[Vertical composition]
Given natural transformations $\alpha: F \Rightarrow G$ and $\beta: G \Rightarrow H$ (all between functors $\cat{C} \to \cat{D}$), their \textbf{vertical composite} $\beta \circ \alpha: F \Rightarrow H$ is defined componentwise:
\[
  (\beta \circ \alpha)_a = \beta_a \circ \alpha_a: Fa \to Ga \to Ha.
\]
\end{definition}

\begin{proposition}
Vertical composition is associative and unital. Hence $[\cat{C},\cat{D}]$ is indeed a category.
\end{proposition}

\begin{proof}
Associativity of vertical composition follows from associativity of composition in $\cat{D}$. The identity natural transformation $\id_F$ (with components $\id_{Fa}$) is the identity for vertical composition. Naturality of the composite is a routine diagram chase.
\end{proof}

\section{Horizontal Composition and Whiskering}

Natural transformations can also be composed in a horizontal direction, combining transformations of adjacent functors.

\begin{definition}[Horizontal composition]
Let $\alpha: F \Rightarrow G$ be a natural transformation between functors $\cat{C} \to \cat{D}$, and $\beta: H \Rightarrow K$ a natural transformation between functors $\cat{D} \to \cat{E}$. The \textbf{horizontal composite} $\beta * \alpha: HF \Rightarrow KG$ has components
\[
  (\beta * \alpha)_a = \beta_{Ga} \circ H(\alpha_a) = K(\alpha_a) \circ \beta_{Fa}.
\]
(Both expressions are equal by naturality of $\beta$.)
\end{definition}

The two-step construction can be visualized as:
\[
  \begin{tikzcd}
    \cat{C} \ar[r, bend left=50, "F", ""{name=F}]
            \ar[r, bend right=50, "G"', ""{name=G}] &
    \cat{D} \ar[r, bend left=50, "H", ""{name=H}]
            \ar[r, bend right=50, "K"', ""{name=K}] &
    \cat{E}
    \ar[Rightarrow, from=F, to=G, "\alpha"]
    \ar[Rightarrow, from=H, to=K, "\beta"]
  \end{tikzcd}
\]

\begin{definition}[Whiskering]
Special cases of horizontal composition:
\begin{itemize}
  \item \textbf{Left whiskering}: for a functor $L: \cat{B} \to \cat{C}$ and a natural transformation $\alpha: F \Rightarrow G$ (between $\cat{C} \to \cat{D}$), define $\alpha L: FL \Rightarrow GL$ with $(\alpha L)_b = \alpha_{Lb}$.
  \item \textbf{Right whiskering}: for $\alpha: F \Rightarrow G$ and a functor $R: \cat{D} \to \cat{E}$, define $R\alpha: RF \Rightarrow RG$ with $(R\alpha)_a = R(\alpha_a)$.
\end{itemize}
\end{definition}

\section{The Interchange Law}

Vertical and horizontal composition interact via the interchange law.

\begin{theorem}[Interchange law]
Let
\[
  \begin{tikzcd}
    \cat{C} \ar[r, bend left=50, "F_1", ""{name=F1}]
            \ar[r, "F_2"', ""{name=F2, above}]
            \ar[r, bend right=50, "F_3"', ""{name=F3}] &
    \cat{D} \ar[r, bend left=50, "G_1", ""{name=G1}]
            \ar[r, "G_2"', ""{name=G2, above}]
            \ar[r, bend right=50, "G_3"', ""{name=G3}] &
    \cat{E}
    \ar[Rightarrow, from=F1, to=F2, "\alpha"]
    \ar[Rightarrow, from=F2, to=F3, "\beta"]
    \ar[Rightarrow, from=G1, to=G2, "\gamma"]
    \ar[Rightarrow, from=G2, to=G3, "\delta"]
  \end{tikzcd}
\]
Then $(\delta \circ \gamma) * (\beta \circ \alpha) = (\delta * \beta) \circ (\gamma * \alpha)$.
\end{theorem}

\begin{proof}
At each object $a \in \cat{C}$, both sides equal $\delta_{F_3 a} \circ G_2(\beta_a) \circ \gamma_{F_2 a} \circ G_1(\alpha_a)$, using naturality of $\delta$ and $\gamma$ at appropriate steps.
\end{proof}

\section{The 2-Category \texorpdfstring{$\Cat$}{Cat}}

The interchange law says that $\Cat$ (or more precisely, $\mathbf{CAT}$, the 2-category of all categories) has the structure of a \textbf{2-category}.

\begin{definition}[2-Category]
A \textbf{2-category} $\mathcal{K}$ consists of:
\begin{itemize}
  \item \textbf{0-cells} (objects): $\cat{C}$, $\cat{D}$, \ldots
  \item \textbf{1-cells} (morphisms): functors $F: \cat{C} \to \cat{D}$.
  \item \textbf{2-cells} (morphisms between morphisms): natural transformations $\alpha: F \Rightarrow G$.
  \item \textbf{Vertical composition} of 2-cells (along shared 1-cells).
  \item \textbf{Horizontal composition} of 2-cells (along shared 0-cells).
  \item The interchange law.
  \item Strict associativity and unit laws at all levels.
\end{itemize}
\end{definition}

$\Cat$ is the archetypal 2-category. Its 0-cells are categories, its 1-cells are functors, and its 2-cells are natural transformations. We will use 2-categorical language sparingly but it provides a useful meta-framework.

\section{Equivalences of Categories Revisited}

We can now make the notion of equivalence precise using natural isomorphisms.

\begin{definition}[Equivalence of categories, revisited]
An \textbf{equivalence of categories} consists of functors $F: \cat{C} \to \cat{D}$ and $G: \cat{D} \to \cat{C}$, together with natural isomorphisms $\eta: \Id_{\cat{C}} \xrightarrow{\sim} GF$ and $\varepsilon: FG \xrightarrow{\sim} \Id_{\cat{D}}$.
\end{definition}

\begin{theorem}
A functor $F: \cat{C} \to \cat{D}$ is part of an equivalence of categories if and only if it is fully faithful and essentially surjective.
\end{theorem}

\begin{proof}[Proof sketch]
($\Rightarrow$) Suppose $F$ is an equivalence via $(F, G, \eta, \varepsilon)$. Full faithfulness: the bijection $\cat{C}(a,b) \cong \cat{D}(Fa, Fb)$ is realized as $f \mapsto Ff$; use $\eta$ and $\varepsilon$ to construct its inverse. Essential surjectivity: for $d \in \cat{D}$, $\varepsilon_d: FG(d) \xrightarrow{\sim} d$ is an isomorphism, so $d \cong F(Gd)$.

($\Leftarrow$) Given $F$ fully faithful and essentially surjective, use essential surjectivity (and the axiom of choice) to select for each $d \in \cat{D}$ an object $Gd$ and an isomorphism $\varepsilon_d: FGd \xrightarrow{\sim} d$. Define $G$ on morphisms using full faithfulness of $F$. Verify $\eta$ and $\varepsilon$ are natural isomorphisms.
\end{proof}

\begin{example}
The skeletal equivalence: every category is equivalent to its \textbf{skeleton} (the full subcategory containing exactly one representative from each isomorphism class of objects).
\end{example}

\section{The Yoneda Lemma, Revisited}

The full strength of the Yoneda lemma is a naturality statement about functor categories.

\begin{theorem}[Yoneda lemma, natural version]
For a locally small category $\cat{C}$, the bijection of Theorem~\ref{thm:yoneda}
\[
  \Nat(\cat{C}(a,-), F) \;\cong\; Fa
\]
is natural in $a \in \cat{C}\op$ and $F \in [\cat{C}, \Set]$. Equivalently, the evaluation functor
\[
  \mathrm{ev}: [\cat{C}, \Set] \times \cat{C} \longrightarrow \Set, \quad (F, a) \mapsto Fa
\]
is naturally isomorphic to the functor $(F, a) \mapsto \Nat(\cat{C}(a,-), F)$.
\end{theorem}

This version makes the Yoneda lemma a statement about a natural isomorphism of bifunctors, and it is the form most useful in practice.

\section{Limits as Representability}
\label{sec:limits-as-rep}

We preview how limits and colimits (Chapter~\ref{ch:limits}) are naturally expressed in terms of representability and natural transformations.

A \textbf{cone} over a diagram $D: \cat{J} \to \cat{C}$ with apex $c$ is a natural transformation $\Delta_c \Rightarrow D$, where $\Delta_c: \cat{J} \to \cat{C}$ is the constant functor at $c$. The \textbf{limit} of $D$, if it exists, is the representing object for the functor
\[
  c \;\mapsto\; \Nat(\Delta_c, D) \;=\; \text{the set of cones over } D \text{ with apex } c.
\]
This makes the definition of limit a clean statement in terms of representable functors.

\section{Exercises}

\begin{exercise}
Verify that vertical composition of natural transformations satisfies the unit and associativity axioms.
\end{exercise}

\begin{exercise}
Show that a natural transformation $\alpha: F \Rightarrow G$ is a natural isomorphism if and only if it has an inverse natural transformation $\alpha^{-1}: G \Rightarrow F$ (with $(\alpha^{-1})_a = (\alpha_a)^{-1}$).
\end{exercise}

\begin{exercise}
Let $F, G: \cat{C} \to \cat{D}$ be naturally isomorphic. Show that for any $H: \cat{D} \to \cat{E}$, we have $HF \cong HG$. Dually, for any $K: \cat{B} \to \cat{C}$, $FK \cong GK$.
\end{exercise}

\begin{exercise}
Show that the functor category $[\mathbf{1}, \cat{D}]$ (where $\mathbf{1}$ is the terminal category) is isomorphic to $\cat{D}$.
\end{exercise}

\begin{exercise}
Show that the functor category $[\mathbf{2}, \cat{C}]$ (where $\mathbf{2}$ is the category $0 \to 1$) is isomorphic to the arrow category $\cat{C}^{\to}$.
\end{exercise}

\begin{exercise}
Prove the interchange law in full detail.
\end{exercise}

\begin{exercise}[$\star$]
Let $\cat{C}$ be a category and $\widehat{\cat{C}} = [\cat{C}\op, \Set]$ the presheaf category. Show that the Yoneda embedding $\yo: \cat{C} \to \widehat{\cat{C}}$ sends limits in $\cat{C}$ to limits in $\widehat{\cat{C}}$. (This requires knowing what limits are; revisit after Chapter~\ref{ch:limits}.)
\end{exercise}

\begin{exercise}[$\star$]
A functor $F: \cat{C} \to \cat{D}$ is an equivalence of categories if and only if there exists a functor $G$ with $GF \cong \Id_{\cat{C}}$ and $FG \cong \Id_{\cat{D}}$. Show that in this case, the same $G$ witnesses $F$ as fully faithful and essentially surjective.
\end{exercise}

\begin{exercise}[$\star\star$]
Prove that for a small category $\cat{C}$, the presheaf category $[\cat{C}\op, \Set]$ is the \textbf{free cocompletion} of $\cat{C}$: every functor $F: \cat{C} \to \cat{D}$ to a cocomplete category $\cat{D}$ extends uniquely (up to natural isomorphism) to a colimit-preserving functor $\bar{F}: [\cat{C}\op,\Set] \to \cat{D}$.
\end{exercise}
